\documentclass[../DoAn.tex]{subfiles}
\begin{document}

Chương này giới thiệu bối cảnh, động lực và định hướng của đề tài. Phần \ref{sec:dvd} phân tích những hạn chế của các cơ chế xác thực hiện tại và đặt ra bài toán nghiên cứu. Phần \ref{sec:hientrang} tổng quan các giải pháp đã có và chỉ ra những khoảng trống còn tồn tại. Phần \ref{sec:muctieu} xác định mục tiêu và định hướng giải pháp của đề tài. Phần \ref{sec:donggop} liệt kê các đóng góp cụ thể. Phần \ref{sec:bocuc} mô tả bố cục của báo cáo.

\section{Đặt vấn đề}
\label{sec:dvd}

Khi thiết bị di động đã trở thành trung tâm lưu trữ thông tin cá nhân và tài chính của rất nhiều người dùng, bài toán xác thực danh tính trên điện thoại thông minh ngày càng nhận được sự quan tâm lớn từ cộng đồng nghiên cứu bảo mật. Trong khi đó, các hình thức tấn công và đánh cắp thông tin ngày càng trở nên tinh vi, đặt ra yêu cầu cấp thiết về các cơ chế bảo vệ hiệu quả hơn.

Các phương pháp xác thực truyền thống đang được sử dụng phổ biến vẫn tồn tại một số những mặt hạn chế. Mã PIN và mật khẩu có thể bị nhìn trộm hoặc bị tấn công vét cạn. Xác thực bằng khuôn mặt có thể bị qua mặt bởi ảnh tĩnh, video hoặc mặt nạ 3D. Vân tay, dù khó tấn công hơn, vẫn có thể bị giả mạo qua dấu để lại trên bề mặt kính và không áp dụng được khi tay người dùng bị ướt hoặc bẩn. Điểm yếu chung và nghiêm trọng nhất của tất cả các hình thức này là chúng chỉ xác thực tại thời điểm mở khóa ban đầu, không kiểm soát liên tục danh tính người dùng trong suốt phiên sử dụng. Việc bổ sung xác thực liên tục theo kiểu truyền thống - yêu cầu nhập lại PIN hoặc quét vân tay định kỳ - lại tạo ra sự đánh đổi khó chấp nhận giữa bảo mật và trải nghiệm người dùng.

Sinh trắc học hành vi (\textit{Behavioral Biometrics}) nổi lên như một hướng tiếp cận đầy hứa hẹn để giải quyết đúng vấn đề trên. Hướng tiếp cận này khai thác các đặc trưng hành vi độc đáo của mỗi cá nhân - cách họ di chuyển, cầm điện thoại, vuốt màn hình, gõ phím - thông qua dữ liệu thu từ các cảm biến đã được tích hợp sẵn trên mọi điện thoại thông minh hiện đại. Ưu điểm nổi bật là khả năng xác thực liên tục và hoàn toàn trong suốt: hệ thống giám sát người dùng trong nền mà không yêu cầu bất kỳ thao tác chủ động nào, qua đó cải thiện trải nghiệm người dùng đồng thời tăng cường bảo mật.

\section{Các giải pháp hiện tại và hạn chế}
\label{sec:hientrang}

Trong những năm gần đây, nhiều hệ thống xác thực hành vi trên thiết bị di động đã được đề xuất và nghiên cứu. IntelliAuth \cite{intelliauth} khai thác dữ liệu cảm biến quán tính kết hợp nhận diện hoạt động để cải thiện độ chính xác trong các ngữ cảnh sử dụng khác nhau. B2auth \cite{b2auth} tập trung vào hành vi touch và được đánh giá trong môi trường triển khai thực tế. M2Auth \cite{m2auth} đề xuất kiến trúc fusion đa phương thức kết hợp nhiều nguồn tín hiệu. MultiLock \cite{multilock} khai thác nhiều lớp xác thực đồng thời nhằm tăng cường độ bền vững của hệ thống.

Tuy nhiên, phần lớn các nghiên cứu hiện có vẫn còn những khoảng trống. Nhiều hệ thống chỉ khai thác một nguồn dữ liệu duy nhất, khiến hiệu quả suy giảm đáng kể trong các tình huống thiết bị đứng yên hoặc người dùng không tương tác trực tiếp với màn hình. Quan trọng hơn, gần như tất cả các hệ thống đều thiếu khả năng phát hiện người dùng lạ (\textit{open-set recognition}) - tức là chủ động từ chối và cảnh báo khi người dùng hiện tại không phải chủ sở hữu thiết bị - đây là điều kiện bắt buộc để một hệ thống có giá trị triển khai thực tế. Ngoài ra, ngay cả khi phát hiện bất thường, các hệ thống hiện tại thường chỉ phản ứng bằng cách khóa màn hình đơn giản, chưa có cơ chế xác thực lại chủ động để khôi phục 
phiên sử dụng một cách an toàn.

\section{Mục tiêu và định hướng giải pháp}
\label{sec:muctieu}

Từ những hạn chế đã phân tích, đề tài hướng tới xây dựng một hệ thống xác thực hành vi liên tục hoàn chỉnh trên nền tảng Android, với ba nhóm mục tiêu cụ thể sau đây.

Về mặt kỹ thuật, đề tài xây dựng một hệ thống xác thực hành vi liên tục chạy hoàn toàn trên thiết bị Android thông qua TensorFlow Lite nhằm đảm bảo quyền riêng tư và khả năng hoạt động ngoại tuyến. Đề tài khảo sát hướng đa phương thức kết hợp dữ liệu cảm biến quán tính 9 kênh với hành vi tương tác màn hình; kết quả đánh giá ở Chương 5 cho thấy nhánh quán tính là đủ và cho hiệu năng tốt hơn ở kịch bản phát hiện người lạ, nên đường quyết định triển khai cuối cùng chỉ sử dụng cảm biến quán tính. Hệ thống tích hợp cơ chế phát hiện người dùng lạ dựa trên ngưỡng điểm tin cậy và cơ chế xác thực dự phòng bằng cử chỉ lắc điện thoại cá nhân khi phát hiện bất thường.

Về phía ứng dụng, đề tài xây dựng hai ứng dụng Android độc lập minh họa toàn bộ luồng hoạt động thực tế: ứng dụng \textit{BioAuth Data Collection} phục vụ thu thập dữ liệu và ứng dụng \textit{BioAuth Authenticator} thực hiện xác thực liên tục trong nền. Ngoài ra, một web demo bằng Streamlit được xây dựng để trực quan hóa quá trình ra quyết định ở mức phiên, trình bày đồng thời điểm của nhánh quán tính, nhánh tương tác và điểm dung hợp cùng hai cột kết luận tương ứng, qua đó minh họa năng lực đa phương thức của hệ thống và cho phép kiểm chứng trên người dùng ngoài tập huấn luyện.

Về mặt nghiên cứu, đề tài thực hiện đánh giá có hệ thống thông qua so sánh hiệu năng giữa các biến thể kiến trúc backbone - CNN 1D, CNN-LSTM và CNN-BiLSTM - theo các chỉ số FAR, FRR và EER, đồng thời đánh giá mức độ đóng góp của từng phương thức và khả năng tổng quát hóa của hệ thống với người dùng chưa xuất hiện trong tập huấn luyện.

Phạm vi của đề tài tập trung vào nền tảng Android phiên bản 10 trở lên với bộ dữ liệu thu từ 26 người tham gia thực hiện ba hoạt động chuẩn hóa là đi bộ, đứng và ngồi; dữ liệu được thu theo giao thức chuẩn hóa trong điều kiện sử dụng tự nhiên và trải trên nhiều dòng thiết bị khác nhau. Đề tài không bao gồm các phương pháp sinh trắc học sinh lý truyền thống, bảo mật tầng mạng, đánh giá chéo có kiểm soát giữa các thiết bị (huấn luyện trên một số dòng máy và kiểm thử trên dòng máy khác), hay cập nhật mô hình trực tuyến. Về mô hình đe dọa, đề tài đánh giá hệ thống trước người lạ \textit{zero-effort} - tức người khác sử dụng thiết bị một cách tự nhiên, không cố bắt chước chủ máy - và không bao gồm việc đánh giá trước các kẻ tấn công chủ động như bắt chước hành vi, phát lại tín hiệu cảm biến, hay quan sát trộm mã cử chỉ lắc.

\section{Đóng góp của đề tài}
\label{sec:donggop}

Đề tài có bốn đóng góp chính sau đây.

Thứ nhất, xây dựng ứng dụng Android thu thập dữ liệu hành vi phục vụ toàn bộ giai đoạn tạo bộ dữ liệu nghiên cứu, từ đó xây dựng được bộ dữ liệu gồm 26 người tham gia với đầy đủ ba loại hoạt động và ba dạng tác vụ tương tác màn hình chuẩn hóa.

Thứ hai, xây dựng pipeline huấn luyện và so sánh có hệ thống ba biến thể kiến trúc backbone cho bộ dữ liệu gồm CNN 1D, CNN-LSTM và CNN-BiLSTM, đồng thời thiết kế hàm quyết định nhận dạng tập mở dựa trên đối sánh embedding với các anchor của chủ sở hữu. Hướng kết hợp đa phương thức với nhánh touch được khảo sát ở vai trò đối chứng nhằm xác định cấu hình phù hợp nhất cho triển khai.

Thứ ba, triển khai toàn bộ hệ thống suy luận trực tiếp trên thiết bị Android thông qua TensorFlow Lite mà không cần kết nối máy chủ, đảm bảo quyền riêng tư người dùng và khả năng hoạt động ngoại tuyến hoàn toàn.

Thứ tư, hoàn tất vòng xác thực bằng một cơ chế dự phòng đơn giản dựa trên cử chỉ lắc điện thoại khi phát hiện người dùng lạ: thay vì quay về mật khẩu hay PIN truyền thống, hệ thống tự khôi phục phiên ngay trong cùng phương thức chuyển động, qua đó khép kín chu trình từ phát hiện bất thường đến khôi phục phiên sử dụng.

\section{Bố cục đồ án}
\label{sec:bocuc}

Phần còn lại của báo cáo được tổ chức thành sáu chương như sau.

Chương 2 trình bày các kiến thức nền tảng cần thiết cho đề tài, bao gồm nguyên lý hoạt động của sinh trắc học hành vi, đặc điểm của dữ liệu cảm biến quán tính và hành vi touch trên Android, các mô hình học máy được sử dụng trong hệ thống, cũng như tổng quan các công trình liên quan và khoảng trống nghiên cứu mà đề tài hướng tới giải quyết.

Chương 3 mô tả quá trình thu thập dữ liệu, bao gồm thiết kế ứng dụng Android thu thập, giao thức thu thập với 26 người tham gia, cùng toàn bộ pipeline tiền xử lý từ dữ liệu thô đến các tensor đầu vào sẵn sàng cho huấn luyện.

Chương 4 trình bày quá trình xây dựng mô hình, gồm huấn luyện và so sánh các biến thể kiến trúc backbone cho dữ liệu IMU, huấn luyện bộ phân loại Random Forest cho dữ liệu touch, thiết kế cơ chế dung hợp đa phương thức và xuất mô hình sang định dạng TFLite.

Chương 5 trình bày thiết kế thực nghiệm và kết quả đánh giá hệ thống theo các chỉ số FAR, FRR và EER, so sánh ba kiến trúc backbone trên hai kịch bản nhận dạng tập đóng và tập mở, phân tích đóng góp của từng phương thức và khả năng tổng quát hóa với người dùng mới, từ đó rút ra cấu hình phù hợp nhất cho triển khai.

Chương 6 mô tả quá trình tích hợp và triển khai cấu hình đã chọn trên Android, bao gồm kiến trúc ứng dụng xác thực, cơ chế suy luận on-device, luồng đăng ký người dùng, dịch vụ xác thực liên tục trong nền và cơ chế xác thực dự phòng bằng cử chỉ lắc.

Chương 7 tổng kết các kết quả đạt được, chỉ ra những hạn chế còn tồn tại và đề xuất các hướng phát triển tiếp theo của hệ thống.

Chương 1 đã phân tích bối cảnh bài toán, chỉ ra khoảng trống của các nghiên cứu hiện tại và xác định định hướng giải pháp của đề tài. Chương 2 tiếp theo sẽ trình bày các nền tảng lý thuyết và tổng quan công trình liên quan làm cơ sở cho toàn bộ hệ thống được xây dựng trong các chương sau.

\end{document}